\cventry{Oct 2020 - Jan 2021}{Product Designer}{\textsc{Adidas}}{Remote}{}
{ Field:  \textbf{Sportswear}
\vspace{0.1cm}
\begin{itemize}
%\item  Brainstorming and conceptualizing innovative ideas for a project. This included sketching, mood boards, and mind mapping to explore various design directions.
%\item Understanding the target audience, industry trends, and competitors is crucial. Analyze market demands and preferences ensure that the design is aligned with consumer expectations.
%\item Specify the physical properties of materials used in the design together with the a color palette selection that reflects the brand or design aesthetics. Tools like Adobe Color or Pantone were used in color selection.
%\item Explanation of general workflow based on the software used (Blender) including refinement, texturing and materials, lighting and rendering
%\item Integrate the 3D design into Unity to create an AR experience. This involved C\# coding  and visual elements for enhancing user experience (UX).
%\item Ensure that the AR design is intuitive, user-friendly, and aligns with the overall brand and project goals.
\item Developed product concepts through sketches, mood boards, and visual research.
\item Researched target users, market trends, and competitors to align design directions with consumer expectations.
\item Defined material properties, color palettes, and visual design guidelines using tools such as Adobe Color and Pantone.
\item Created and refined 3D assets in Blender, including materials, lighting, texturing, and rendering.
\item Integrated 3D assets into Unity to build an AR experience with C# interactions.
\item Improved the AR user experience by aligning visual, interaction, and brand requirements.
\end{itemize}}